%! AP Statistics — Unit 1, Lesson 1.9 — Homework
\documentclass[10pt]{article}
\usepackage{apstats-article}
\usepackage{apstats-boxes}

%=============================================================================
\begin{document}

\pageheader{Unit 1, Lesson 1.9}{Homework: Comparing Distributions}
\namedateperiod

\vspace{0.10in}

\begin{tcolorbox}[colback=sky, colframe=navy, boxrule=0.8pt, arc=2mm,
    left=3mm, right=3mm, top=1.5mm, bottom=1.5mm]
\small Use explicit comparative language in every sentence that compares groups.
Interpret statistics in context. Show all fence calculations.
\end{tcolorbox}

\vspace{0.12in}

% ════════════════════════════════════════════════════════════════════════════
% PART A — Read and compare parallel boxplots
% ════════════════════════════════════════════════════════════════════════════
\begin{blurbbox}[Part A \quad Read and Compare Parallel Boxplots]{navylight}
\small

A school principal recorded the number of \textbf{books read per semester}
by students in a morning class and an afternoon class.

\vspace{0.10in}
\begin{center}
\begin{tikzpicture}[scale=1.0]
  % Axis: 50 to 95, x = (v-50)*0.16, label every 5
  \draw[thick] (0.2, 0) -- (7.40, 0);
  \foreach \v/\x in {50/0.00, 55/0.80, 60/1.60, 65/2.40, 70/3.20,
                      75/4.00, 80/4.80, 85/5.60, 90/6.40, 95/7.20} {
    \draw (\x, -0.12) -- (\x, 0.12);
    \node[below, font=\small] at (\x, -0.12) {\v};
  }
  \node[below, font=\small\itshape] at (3.60, -0.56) {Books read per semester};

  % Morning: min=60(1.60), Q1=68(2.88), M=73(3.68), Q3=79(4.64), max=90(6.40)
  % IQR=11, upper fence=79+16.5=95.5. No outliers.
  % y_center=1.10
  \node[left, font=\small] at (0.0, 1.10) {Morning};
  \draw[thick, navy] (1.60, 1.10) -- (2.88, 1.10);
  \draw[thick, navy] (2.88, 0.85) rectangle (4.64, 1.35);
  \draw[thick, navy] (3.68, 0.85) -- (3.68, 1.35);
  \draw[thick, navy] (4.64, 1.10) -- (6.40, 1.10);
  \draw[thick, navy] (1.60, 0.95) -- (1.60, 1.25);
  \draw[thick, navy] (6.40, 0.95) -- (6.40, 1.25);

  % Afternoon: min=55(0.80), Q1=62(1.92), M=70(3.20), Q3=78(4.48), max=88(6.08)
  % IQR=16, upper fence=78+24=102. No outliers.
  % y_center=0.35
  \node[left, font=\small] at (0.0, 0.35) {Afternoon};
  \draw[thick, navy] (0.80, 0.35) -- (1.92, 0.35);
  \draw[thick, navy] (1.92, 0.10) rectangle (4.48, 0.60);
  \draw[thick, navy] (3.20, 0.10) -- (3.20, 0.60);
  \draw[thick, navy] (4.48, 0.35) -- (6.08, 0.35);
  \draw[thick, navy] (0.80, 0.20) -- (0.80, 0.50);
  \draw[thick, navy] (6.08, 0.20) -- (6.08, 0.50);
\end{tikzpicture}
\end{center}

\vspace{0.10in}

\begin{enumerate}[leftmargin=1.5em, itemsep=8pt]

  \item Read the five-number summary for each class.\\[5pt]
        \renewcommand{\arraystretch}{1.5}
        \begin{tabularx}{\linewidth}{@{} l Y Y Y Y Y @{}}
        \rowcolor{sky}
        & $\min$ & $Q_1$ & $M$ & $Q_3$ & $\max$ \\
        \midrule
        Morning   & \blank{1.2cm} & \blank{1.2cm} & \blank{1.2cm} & \blank{1.2cm} & \blank{1.2cm} \\
        Afternoon & \blank{1.2cm} & \blank{1.2cm} & \blank{1.2cm} & \blank{1.2cm} & \blank{1.2cm} \\
        \end{tabularx}\\[4pt]
        Morning IQR $= \blank{2cm}$ books \hfill
        Afternoon IQR $= \blank{2cm}$ books

  \item Write one AP-quality sentence comparing the \textbf{centers} of the
        two distributions. Name both classes, use a comparison word, give
        both medians, and include context.\\[5pt]
        \writeline\\[1pt]\writeline

  \item Write one AP-quality sentence comparing the \textbf{spreads}.\\[5pt]
        \writeline\\[1pt]\writeline

  \item Which class, if either, shows more overlap between the two
        distributions? What does a large overlap suggest about the groups?\\[5pt]
        \writeline\\[1pt]\writeline

\end{enumerate}

\end{blurbbox}

\vspace{0.12in}

% ════════════════════════════════════════════════════════════════════════════
% PART B — Back-to-back stemplot
% ════════════════════════════════════════════════════════════════════════════
\begin{blurbbox}[Part B \quad Back-to-Back Stemplot]{greenacc}
\small

The sorted quiz scores (out of 100) for two groups are listed below.\\[4pt]
\textbf{Group X} ($n = 8$):\enspace 62, 65, 67, 71, 74, 76, 81, 83\\
\textbf{Group Y} ($n = 8$):\enspace 64, 68, 72, 75, 78, 82, 85, 89

\vspace{0.10in}

\begin{enumerate}[leftmargin=1.5em, itemsep=8pt]

  \item Fill in the missing leaves in the back-to-back stemplot.
        Group X leaves are written in \textit{reverse} order on the left.\\[6pt]
        \begin{center}
        \small
        \renewcommand{\arraystretch}{1.4}
        \begin{tabular}{r c l}
        \multicolumn{1}{c}{\textit{Group X (reversed)}} & Stem & \textit{Group Y} \\
        \hline
        \blank{2.5cm} & 6 & \blank{2.5cm} \\
        \blank{2.5cm} & 7 & \blank{2.5cm} \\
        \blank{2.5cm} & 8 & \blank{2.5cm} \\
        \hline
        \multicolumn{3}{c}{\small Key: $7 \mid 4 = 74$ for Group X;\enspace $7 \mid 5 = 75$ for Group Y}
        \end{tabular}
        \end{center}

  \item Find the five-number summary for each group from the data.\\[6pt]
        \renewcommand{\arraystretch}{1.5}
        \begin{tabularx}{\linewidth}{@{} l Y Y Y Y Y @{}}
        \rowcolor{sky}
        & $\min$ & $Q_1$ & $M$ & $Q_3$ & $\max$ \\
        \midrule
        Group X & \blank{1.2cm} & \blank{1.2cm} & \blank{1.2cm} & \blank{1.2cm} & \blank{1.2cm} \\
        Group Y & \blank{1.2cm} & \blank{1.2cm} & \blank{1.2cm} & \blank{1.2cm} & \blank{1.2cm} \\
        \end{tabularx}

  \item Using the stemplot and the five-number summaries, write a complete
        AP comparison of the two groups' quiz score distributions.
        Address shape, center, and spread.\\[5pt]
        \writeline\\[1pt]\writeline\\[1pt]\writeline\\[1pt]\writeline

\end{enumerate}

\end{blurbbox}

\vspace{0.12in}

% ════════════════════════════════════════════════════════════════════════════
% PART C — AP-Style Free Response
% ════════════════════════════════════════════════════════════════════════════
\begin{blurbbox}[Part C \quad AP-Style Free Response]{redacc}
\small

Using the parallel boxplots from Part A, write a \textbf{complete AP comparison}
of the books-read distributions for the morning and afternoon classes. Address
all four SOCS components. Every sentence must use explicit comparative language
that names both classes.

\vspace{0.08in}
\writeline\\[1pt]\writeline\\[1pt]\writeline\\[1pt]
\writeline\\[1pt]\writeline\\[1pt]\writeline

\end{blurbbox}

\end{document}
